%% ---------------------------------------------------------------------
\section{Reflected triangles and a pumping lemma}\label{sec:pumping}

\subsection{Relative coordinates}

Fix $n\ge2$. For a cell $(t,i)$ of the line of length $n$ put
$\tau:=t-(n-1)$ (time elapsed since the front reached the right end) and
$\kappa:=n-i$ (distance to the right end). Then
\[
R_n=\{(\tau,\kappa):\ 0\le\kappa\le\tau\le n-1\},
\]
the right border sits at $\kappa=-1$ and the left border at $\kappa=n$.
A cell $(\tau,\kappa)\in R_n$ is computed from
$(\tau-1,\kappa+1),(\tau-1,\kappa),(\tau-1,\kappa-1)$; each of these is
either in $R_n$, or a border, or satisfies $\tau'-\kappa'\in\{-1,-2\}$,
i.e.\ lies on one of the anti-diagonals $t+i=2n-2$, $t+i=2n-3$, where
$C_n$ coincides with $C_\infty$ by \cref{lem:agree}. We call
\[
J_n(\kappa):=\bigl(C_\infty(n-3+\kappa,\,n-\kappa),\;C_\infty(n-2+\kappa,\,n-\kappa)\bigr),
\qquad 0\le\kappa\le n-1,
\]
the \emph{input word} of $R_n$; its two components are the \emph{lower}
($t+i=2n-3$) and \emph{upper} ($t+i=2n-2$) input cells.

\begin{lemma}[determinacy]\label{lem:determinacy}
There is a single partial map $\Psi$, depending on $\delta$ only, such that
for every $n\ge2$ and every $(\tau,\kappa)\in R_n$,
$C_n\bigl(n-1+\tau,\,n-\kappa\bigr)$ is obtained by evaluating $\Psi$ at
$(\tau,\kappa)$ on the input word $J_n$ with borders at $\kappa=-1$ and
$\kappa=n$. In particular the value only depends on the input cells and
borders in the backward dependency cone of $(\tau,\kappa)$.
\end{lemma}

\begin{proof}
The recurrence in relative coordinates does not mention $n$; only the
position of the left border does.
\end{proof}

For $(\tau,\kappa)\in R_n$ let $\mathrm{cone}_n(\tau,\kappa)$ be the set of
input cells (pairs $(\kappa',\text{lower/upper})$) reached from
$(\tau,\kappa)$ by repeatedly passing to the three predecessors, stopping at
input cells and borders; say that the cone is \emph{right-anchored} if it
does not reach the left border $\kappa'=n$.

\begin{lemma}\label{lem:cone}
For $n\ge4$ the cone of the right end at the firing time,
$\mathrm{cone}_n(n-1,0)$, is right-anchored and consists of the upper input
cell at $\kappa'=0$, both input cells at $1\le\kappa'\le\lfloor n/2\rfloor$,
and, when $n$ is odd, the lower input cell at $\kappa'=(n+1)/2$.
\end{lemma}

\begin{proof}
A predecessor step changes $\tau$ by $-1$ and $\kappa$ by at most one, and
input cells satisfy $\tau'=\kappa'-1$ or $\kappa'-2$. Starting from
$(n-1,0)$, a path that reaches an input cell at position $\kappa'$ after $s$
steps satisfies $\kappa'\le s$ and $n-1-s=\kappa'-1$ or $\kappa'-2$; hence
$2\kappa'\le n+1$ (upper cells: $2\kappa'\le n$). Conversely every such cell
is reached by the path that first moves left $\kappa'$ times and then stays
put (the intermediate cells lie in $R_n$). The largest $\kappa'$ reached is
$\lceil n/2\rceil\le n-1$ for $n\ge 4$, so the left border is never reached;
the right border only contributes the constant symbol~$\ast$. (The same set
is also computed mechanically in the accompanying code, file
\texttt{pumping.py}.)
\end{proof}

\begin{lemma}[pumping]\label{lem:pumping}
Let $\delta$ be a minimal-time solution and $4\le n<n'$. Then $J_n$ and
$J_{n'}$ differ on $\mathrm{cone}_n(n-1,0)$.
\end{lemma}

\begin{proof}
Suppose they agree there. By \cref{lem:determinacy,lem:cone} the relative
cell $(n-1,0)$ of $R_{n'}$ (which belongs to $R_{n'}$ since $n-1\le n'-1$,
and whose cone is the same set of input positions, far from the left border
$\kappa=n'>n$) has the same value as the relative cell $(n-1,0)$ of $R_n$,
namely $C_n(2n-2,n)=\mathsf F$. Thus $C_{n'}(n'+n-2,n')=\mathsf F$ although
$n'+n-2<2n'-2$, contradicting minimality for $n'$.
\end{proof}

The point of \cref{lem:pumping} is that it constrains the \emph{half-line
diagram only}: no reflected triangle has to be computed. It is the source of
the constraint family $\mathrm{PUMP}(n,n')$ used in \cref{sec:encoding}. As a
sanity check, we verified mechanically that Mazoyer's six-state solution
satisfies $\mathrm{PUMP}(n,n')$ for all $4\le n<n'\le120$.

\subsection{Consequence: a speed-$1/3$ barrier}

Write the half-line diagram in the frame of the front:
$\Phi_t(j):=C_\infty(t,\,t+1-j)$ for $0\le j\le t$ (the cell at
\emph{depth}~$j$ behind the front at time~$t$). By the recurrence,
\begin{equation}\label{eq:comoving}
\Phi_{t+1}(j)=\delta\bigl(\Phi_t(j),\,\Phi_t(j-1),\,\Phi_t(j-2)\bigr)
\qquad(0\le j\le t),
\end{equation}
with $\Phi_t(-1)=\Phi_t(-2)=\mathsf L$: in this frame the half-line is a
\emph{one-sided} automaton in which information only travels to larger
depths, at most two depths per step. In particular every depth-row
$t\mapsto\Phi_t(j)$ is eventually periodic.
In these coordinates
$J_n(\kappa)=\bigl(\Phi_{n-3+\kappa}(2\kappa-2),\ \Phi_{n-2+\kappa}(2\kappa-1)\bigr)$.

\begin{theorem}[speed-$1/3$ barrier]\label{thm:third}
Let $\delta$ be a minimal-time solution and suppose that the depth-rows have
a common eventual period: there are $P\ge1$ and thresholds $T(j)$ with
$\Phi_{t+P}(j)=\Phi_t(j)$ for all $t\ge T(j)$. Then for every $n\ge4$ there
is a depth $j\le n$ with $T(j)>n+\tfrac{j}{2}-\tfrac52$. Consequently
$\limsup_{j\to\infty}T(j)/j\ge 3/2$; in the original coordinates, the
periodic regime behind the front cannot overtake the line $i=t/3$.
\end{theorem}

\begin{proof}
If every input cell of $\mathrm{cone}_n(n-1,0)$ is read in the periodic
regime of its row, i.e.\ $n-3+\kappa\ge T(2\kappa-2)$ and
$n-2+\kappa\ge T(2\kappa-1)$ for all cone positions, then
$J_{n+P}=J_n$ on the cone, contradicting \cref{lem:pumping}. The cone
positions satisfy $\kappa\le\lceil n/2\rceil$, so the violated inequality
concerns a depth $j\in\{2\kappa-2,2\kappa-1\}$, $j\le n$, with
$T(j)>n-3+\kappa\ge n+\frac{j}{2}-\frac52$. If $T(j)\le\alpha j+\beta$ for
all~$j$ with $\alpha<3/2$, then for $n$ large and all $j\le n$ we have
$\alpha j+\beta\le n+\frac j2-\frac52$ (as $(\alpha-\frac12)j\le(\alpha-\frac12)n
<n-\beta-\frac52$), a contradiction. A cell of depth $j$ at time $T(j)$ has
position $T(j)+1-j$, whence the geometric statement.
\end{proof}

\Cref{thm:third} explains a posteriori the ubiquity of the speed-$1/3$
signal emitted by the general in the classical constructions
\cite{Waksman66,Balzer67,Mazoyer87}: its crossing with the anti-diagonal
$t+i=2n-2$ is exactly the far end of $\mathrm{cone}_n(n-1,0)$. It also
explains the failure mode of the ``lucky'' partial solutions found by the
solvers, whose half-line diagrams become periodic behind the front at
speed larger than $1/3$ (see \cref{sec:experiments}).
